\section{Choosing the dual forest using raw cancellation obstructions}
\label{sec:forests}

\subsection{The restricted certificate class}
Fix a zero-/one-handle spanning forest $F$ in the core.  Contracting $F$
leaves one zero-handle in each component and some one-handles.  Every
two-handle has a geometric attaching cycle through their belt disks.  Let
\[
          A=(a_{rc})\in\N^{R\times C}
\]
be its unsigned occurrence matrix: $r$ labels a two-handle and $c$ a
surviving one-handle.  Signs and cancellations in the fundamental group are
not used.  Multiple geometric crossings contribute multiple occurrences.
One may saturate each count at two because only zero, one, and at least two
matter to the predicates below.

The same row set $R$ labels the edges of the rooted dual graph $D$.
When interface relations are included, give each row ineligible for top
cancellation an artificial root loop.  Such a row cannot be selected in
a tree; it still participates fully in every stopping test.
Choosing a dual spanning tree $T\subseteq R$ deletes its rows, by
Section~\ref{sec:geometry}.  We ask whether all remaining rows can be
removed by raw geometric one-/two-handle cancellation.  A raw pivot
$(r,c)$ requires $a_{rc}=1$ and has the unsigned update
\begin{equation}\label{eq:pivot}
       a'_{ih}=a_{ih}+a_{ic}a_{rh}
       \qquad (i\ne r,\ h\ne c),
\end{equation}
followed by deletion of $r$ and $c$.  This counts the crossings introduced
by pushing the other attaching bands through the cancelling pair.  It is
the conservative raw update used by the existing geometric certifier.
Saturating after each update preserves every future zero/one test.

\begin{definition}[Forest-compatible raw completion]
The pair $(D,A)$ is \emph{completable} if some spanning tree $T$ of $D$
leaves $A[R\setminus T,C]$ completely eliminable by pivots
\eqref{eq:pivot}.  The lower forest $F$ remains fixed throughout this
definition.
\end{definition}
This problem deliberately excludes normal isotopies that change occurrence
counts, arbitrary handle slides, and replacement of the lower forest.
Its positive answers are useful geometric certificates.  Its negative
answers have only the scope of this definition.

\subsection{Stopping sets and the exact raw criterion}
We recall the stopping-set theorem of the incoming stopping-core work
\cite{StoppingCores,SaturatedBundles}, including a proof because its invariant
is essential to the new search.

\begin{definition}[Stopping set]
A nonempty set $S\subseteq R$ is a stopping set when
\begin{equation}\label{eq:stopping}
          \sum_{r\in S}a_{rc}\ne1
          \quad\hbox{for every column }c.
\end{equation}
Equivalently every column has zero or at least two occurrences in $S$.
A \emph{unit-column peel} removes the unique row incident with a column
whose total is one.  The column can be removed as well: it has no other
occurrences, so this pivot creates no fill.
\end{definition}

\begin{lemma}[Canonical stopping residual]\label{lem:peeling}
Exhaustive unit-column peeling leaves the same row set regardless of
tie-breaking.  This set is the union of all stopping sets.  In particular
a row restriction is peelable exactly when it contains no stopping set,
and peelability is hereditary under deleting rows.
\end{lemma}
\begin{proof}
The union of stopping sets is stopping.  If a column had total one on the
union, its unique contributing row would belong to a constituent stopping
set on which that column also had total one, a contradiction.

A row in a stopping set cannot be peeled while the stopping set is intact,
because every column through that row has at least two occurrences within
the set.  Thus every stopping set survives every sequence of peels.  At
termination, if rows remain, their set is itself stopping.  It therefore
equals the union of all stopping sets and is independent of the order.
Finally a stopping subset of a smaller row restriction would also be a
stopping subset of the larger restriction, proving heredity.
\end{proof}

\begin{lemma}[Completeness for raw pivots]\label{lem:raw-complete}
All rows can be removed by arbitrary raw pivots \eqref{eq:pivot} if and only
if unit-column peeling removes all rows.
\end{lemma}
\begin{proof}
Peeling is a particular raw pivot sequence.  For the converse, suppose
$S$ is stopping and perform a raw pivot $(r,c)$.  If $r\notin S$, the
total $m=\sum_{i\in S}a_{ic}$ is zero or at least two.  In every remaining
column $h$, the new total is
\[
       \sum_{i\in S}a_{ih}+m a_{rh}.
\]
It is zero or at least two: a formerly zero total either stays zero or
increases by at least two, and a positive old total was already at least
two.  Thus $S$ remains stopping after removal of $c$.

If $r\in S$, then $m\geq2$ because $a_{rc}=1$.  Put
$S'=S\setminus\{r\}$.  The set $S'$ is nonempty; otherwise $S$ would
already have total one in column $c$.  In a remaining column its total is
\[
 \sum_{i\in S'}a_{ih}+(m-1)a_{rh}
       \ \geq\ \sum_{i\in S}a_{ih}.
\]
An old positive total was at least two, and an old zero total has all its
summands zero.  Hence $S'$ is stopping.  Inductively a nonempty stopping
set survives every pivot sequence, preventing complete elimination.
Lemma~\ref{lem:peeling} now proves the converse.
\end{proof}

The proof explains why signed exponent sums are inadequate.  A curve with
raw word $x x^{-1}x$ crosses the belt disk three times, even though its
group word freely reduces to $x$.  Treating these as one crossing silently
changes the permitted geometric move class.

\subsection{Spanning trees as transversals of obstructions}
\begin{theorem}[Forest--stopping-set equivalence]\label{thm:transversal}
Assume $D$ is connected after boundary roots are identified.  A spanning
tree $T$ is a successful forest-compatible raw completion if and only if
\begin{equation}\label{eq:hitting}
                   T\cap S\ne\varnothing
       \qquad\hbox{for every stopping set }S\subseteq R.
\end{equation}
Thus the problem is to choose a basis of a graphic matroid that is a
transversal of the stopping-set family.
\end{theorem}
\begin{proof}
If $T$ misses $S$, the same stopping set survives in $R\setminus T$, so
Lemma~\ref{lem:raw-complete} forbids completion.  Conversely, if the
restriction to $R\setminus T$ fails to peel, its nonempty canonical
residual is a stopping set missed by $T$.  Taking the contrapositive proves
the other direction.
\end{proof}

The family in \eqref{eq:hitting} can be exponentially large.  The theorem
is useful because a failed peel exposes a witness without listing the
whole family.  Moreover, peeling must not irrevocably reserve a row for
lower cancellation.  A row peeled during a partial search may still be
chosen later as a dual-tree edge.  Heredity makes this safe.

\subsection{Complete obstruction-guided search}
Maintain a selected forest $Q$ of dual edges.  Contract it in $D$ and remove
its rows from $A$.  Compute the stopping residual $K_Q$.  If it is empty,
extend $Q$ to any spanning tree of $D$ and finish.  Otherwise choose any
nonempty stopping subset $S\subseteq K_Q$, and branch on the edges
$e\in S$ that are not loops in $D/Q$.  The child adds $e$ to $Q$.
An inclusion-minimal stopping witness can be sought by successive row
removals and re-peeling, but minimality is not required for correctness.

\begin{theorem}[Complete forest search]\label{thm:search}
The preceding search is sound and complete for forest-compatible raw
completion.  Let $s=|V(D)|-1$ and let $w$ bound the sizes of the stopping
witnesses actually used.  Its branching depth is at most $s$, with at most
\begin{equation}\label{eq:search-nodes}
                     1+w+w^2+\cdots+w^s
\end{equation}
nodes.  If a node, including witness selection, costs $P(N)$ for encoded
input size $N$, the total cost is bounded by \eqref{eq:search-nodes} times
$P(N)$.  In particular it is $w^{O(s)}\poly(N)$ when $w\geq2$ and the
witness routine is polynomial; for $w=1$ it is $O((s+1)P(N))$.
\end{theorem}
\begin{proof}
A successful spanning-tree extension of $Q$ must meet any stopping witness
$S$ in a row not already in $Q$.  Every such edge is independent of $Q$,
so it appears in one of the stated branches.  Conversely each branch
retains independence and deletes exactly the rows that its chosen dual
edges cancel.  Thus branching neither omits a possible completion nor
creates a false one.

If $K_Q=\varnothing$, connectivity of $D$ ensures that the forest $Q$
extends to a spanning tree $T$.  The row set $R\setminus T$
is a subset of the already peelable set $R\setminus Q$.  Heredity proves
completion.  The final peel order should be recomputed or checked for
this actual $T$, rather than treating the earlier exploratory order as
an unverified certificate.

Each branch increases the forest rank by one.  A forest on $s+1$ vertices
has at most $s$ edges, proving the depth and node bounds.
\end{proof}

\begin{corollary}[A local quasi-polynomial regime]
If the source size and coordinate bit lengths are polynomial in a diagram
size $n$, and $s\log(w+1)=O((\log n)^k)$ for fixed $k$, this restricted
forest search has quasi-polynomial bit complexity.  The polynomial
degree-two-column algorithm in Section~\ref{sec:sparse} requires no small
bound on $s$.
\end{corollary}
This is a local certificate-class bound.  The source triangulation can have
linearly many interior vertices, and no bound on the total hierarchy work
follows from it alone.

\begin{proposition}[Packing obstruction]\label{prop:packing}
Let $Q$ be the currently selected forest, and let
$S_1,\ldots,S_p$ be pairwise disjoint nonempty stopping sets in
$R\setminus Q$.  A necessary condition for completion is
\begin{equation}\label{eq:packing}
       \rk_{D/Q}\left(\bigcup_{i=1}^p S_i\right)\geq p.
\end{equation}
For an edge set $X$, the rank is the number of vertices minus the number
of connected components of the graph with edge set $X$, including isolated
vertices and ignoring loops as rank contributions.
\end{proposition}
\begin{proof}
A completion must choose at least one edge from each $S_i$.  Disjointness
forces at least $p$ distinct chosen edges in their union.  These are part
of an independent forest in the contracted graph, so its rank is at
least $p$.
\end{proof}
In particular a stopping singleton is a mandatory tree edge.  If mandatory
edges form a cycle, no forest-compatible completion exists.  A stopping
witness whose every edge becomes a loop is the rank-zero case.  These
certificates use exact graph incidences and exact integer tests.

\subsection{Why tree choice matters, even for a ball}
\begin{example}[A geometric two-edge choice]\label{ex:choice}
Start with a solid torus, equipped with a fixed meridian belt disk for its
one-handle.  Choose two disjoint parallel primitive curves on its boundary.
Arrange the first to cross the belt disk once.  In the annulus complementary
to the first curve, a finger isotopy of the second adds a cancelling pair
of crossings, so the second crosses three times while remaining embedded
and parallel to the first.  Attach two two-handles along the curves and cap
one of the two resulting sphere boundaries with a three-handle.  Before
the cap the manifold is $S^2\times I$; after it the manifold is a ball.

The dual graph has one nonroot vertex and two parallel edges $a,b$ to the
root.  The raw occurrence matrix is
\[
                  A=\begin{pmatrix}1\\3\end{pmatrix}.
\]
Deleting $b$ leaves a raw singleton and succeeds.  Deleting $a$ leaves a
stopping singleton and fails.  Both choices are valid topological
two-/three-handle cancellations.  The difference lies entirely in the
subsequent restricted raw certificate.

If both primitive curves are given such finger pairs, the matrix becomes
$(3,3)^{\mathsf T}$.  Each row is a stopping singleton.  Their dual edges
have rank one, so \eqref{eq:packing} certifies impossibility in the raw
forest class.  The manifold is still a ball.  Thus even an exhaustive
negative answer in this class must remain inconclusive for handlebody
recognition.
\end{example}

This is a geometric existence example in a general handle structure.
Its actual column totals are four and six, respectively, so it is not
a native triangular-face core with the specified cocore.  Capping counts
preserves the finite cancellation decision, but does not authenticate a
native realization.  The example establishes that forest choice can
matter geometrically; the source tests provide the separate native evidence.

\begin{example}[Peelable row sets are not a matroid]\label{ex:nonmatroid}
Consider rows
\[
 r_0=(0,0,1),\quad r_1=(0,1,1),\quad
 r_2=(1,1,0),\quad r_3=(1,1,1).
\]
The row sets $X=\{r_2,r_3\}$ and $Y=\{r_0,r_1,r_2\}$ both peel and
$|X|<|Y|$.  But the column totals of $X\cup\{r_0\}$ are $(2,2,2)$,
and those of $X\cup\{r_1\}$ are $(2,3,2)$.  Both enlarged sets are
stopping.  The augmentation axiom therefore fails.  One cannot solve the
general problem merely by declaring peelable row sets to be a second
matroid.  The next section obtains a matroid formulation from a genuine
structural restriction on column support.
\end{example}
